\documentclass[10pt,a4paper]{amsart}
\usepackage[margin=19mm]{geometry}
\usepackage{amsmath,amssymb,mathtools}
\usepackage[scr=rsfs,cal=euler]{mathalfa}
\usepackage{xfrac}
\usepackage[unicode,hidelinks,bookmarks=false]{hyperref}
\newcommand{\ie}{i.e.\ }
\newcommand{\cf}{cf.\ }
\newcommand{\resp}{respectively\ }
\newcommand{\Subsection}{\subsection}
\providecommand{\nobreakdash}{\nobreak\hbox{-}}
\setlength{\emergencystretch}{2em}
\multlinegap0pt
\usepackage{bm}
\usepackage{bbm}
\usepackage{dsfont}
\usepackage{xspace}
\usepackage{cancel}
\usepackage{mathtools}
\usepackage{amsthm}
\makeatletter
\@ifundefined{theorem}{\newtheorem{theorem}{Theorem}}{}
\@ifundefined{proposition}{\newtheorem{proposition}[theorem]{Proposition}}{}
\makeatother
\newcommand{\BV}{\operatorname{BV}}
\newcommand{\Borel}{\mathscr{B}}
\newcommand{\Circ}{\operatorname{Circ}}
\newcommand{\C}{{\mathbb C}}
\newcommand{\R}{{\mathbb R}}
\newcommand{\Z}{{\mathbb Z}}
\newcommand{\e}{\ensuremath{\mkern.25\thinmuskip\mathrm{e}\mkern.4\thinmuskip}}
\newcommand{\im}{\operatorname{Im}}
\newcommand{\mfrac}[2]{\tfrac{#1\strut}{#2\strut}}
\newcommand{\orbit}{\mathbb{A}}
\newcommand{\pms}{{\scalebox{0.7}{$\pmb{\pm}$}}}
\newcommand{\sminus}{\setminus} %\newcommand{\sminus}{\!\smallsetminus\!}
\renewcommand{\Re}{\operatorname{Re}}
\renewcommand{\d}{\ensuremath{\operatorname{d}\!}}
\renewcommand{\i}{\mathrm{i}\mkern.25\thinmuskip}
\title[Reading analytic invariants of parabolic diffeomorphisms fro]{Reading analytic invariants of parabolic diffeomorphisms from their orbits}
\author{Martin Klimes, Pavao Mardesic, Goran Radunovic, Maja Resman}
\date{Source: arXiv:2112.14324v2 (CC BY 4.0)}
\begin{document}
\maketitle

\noindent\emph{Source and licence.} Reading analytic invariants of parabolic diffeomorphisms from their orbits. Version arXiv:2112.14324v2 (CC BY 4.0), distributed under the Creative Commons Attribution 4.0 International licence, as recorded in the arXiv licence metadata for this exact version and shown on the abstract page https://arxiv.org/abs/2112.14324v2. Full rights notice: licenses/arxiv-2112.14324-CC-BY-4.0.txt.\par\smallskip

\noindent\textit{Editorial scope.} Counted unit: the Proposition whose source label is \texttt{prop:druga} in the article (source0.tex lines 1242--1250, source label \texttt{prop:druga}), delivered together with the article's own complete original proof of it (source0.tex lines 1252--1306, one \texttt{proof} environment of 55 lines). No mathematics is written, repaired or re-derived: statement and proof are the article's own text line for line. Required context retained uncounted: eq:thetaalpha (lines 1009--1011); eq:ThetaC (lines 1092--1094); eq:Q (lines 1231--1236). Every definition, display and label of those lines is retained unchanged. Omitted from this package and not redistributed: 1--1008, 1012--1091, 1095--1230, 1237--1241, 1251--1251, 1307--1465. Nothing inside a retained excerpt is omitted or altered: every retained body is delivered exactly as the article's lines are written, so no omission receipt of any kind accompanies this package. Citations: the retained excerpts carry no citation command at all, so no bibliography entry of the article is transcribed and no reference is redistributed. Reference adaptation: none was needed and none was made; every reference inside the delivered unit resolves inside it, so no unresolved reference or citation is produced. Printed slips kept literal: no printed word, symbol, hypothesis or number of the counted statement or of its proof was altered. Layout and class: the article's own class is \texttt{amsart} with packages including inputenc, fontenc, amsmath, amsthm, amssymb, amsfonts, mathtools, bbm, stmaryrd, mathrsfs, caption, subcaption, enumitem, tikz; the wrapper is the lane's pinned \texttt{amsart} editorial wrapper at 10pt on A4, which restates the theorem-like environments the retained text needs and adds the article's own macro definitions that the retained text needs (\texttt{\textbackslash BV}, \texttt{\textbackslash Borel}, \texttt{\textbackslash C}, \texttt{\textbackslash Circ}, \texttt{\textbackslash R}, \texttt{\textbackslash Re}, \texttt{\textbackslash Z}, \texttt{\textbackslash d}, \texttt{\textbackslash e}, \texttt{\textbackslash i}, \texttt{\textbackslash im}, \texttt{\textbackslash mfrac}, \texttt{\textbackslash orbit}, \texttt{\textbackslash pms}, \texttt{\textbackslash sminus}), with extra wrapper packages (bm, bbm, dsfont, xspace, cancel, mathtools). The delivered numbering is the unit's own. Assets: no retained span contains a figure environment, a graphics-inclusion command, a PDF-page inclusion, a table or a TikZ picture; no image file is copied into this package, the package declares no asset dependency and no image file is redistributed with it. Licence: version arXiv:2112.14324v2 of this article is distributed under the Creative Commons Attribution 4.0 International licence, as recorded in the arXiv licence metadata for this exact version and shown on the abstract page https://arxiv.org/abs/2112.14324v2. A full rights notice is delivered at licenses/arxiv-2112.14324-CC-BY-4.0.txt. Characters: every retained excerpt is pure ASCII, so the delivered engine, XeLaTeX with its default Latin Modern text font, reports no missing character.
\begin{equation}\label{eq:thetaalpha}
	\alpha\in\ ]\,\underline{\theta},\overline{\theta}\,[\ = \ ]\tfrac{\pi}{2},\tfrac{3\pi}{2}[\mod 2\pi\Z,
\end{equation} 

\begin{equation}\label{eq:ThetaC}
\Theta_\orbit(s)=\int_{b+\i \e^{-\i\alpha}\R}\!\!(B_m\!\circ\psi_\orbit(x))\,\e^{-s\,t(x)}\d\psi_\orbit(x),\qquad b\in\ ]\!-\!1,0[.
\end{equation}

\begin{equation}\label{eq:Q}
\begin{aligned}
&\tilde Q(s,\sigma):=\widetilde\Borel_{\alpha}\{\e^{-\sigma u(x)}\}(s),\\
&Q(s,\sigma):=\Borel_{\alpha}\{\e^{-\sigma u(x)}\}(s)=\BV[\tilde Q(s,\sigma)]_{s\in \e^{\i\alpha}\R_{\geq 0}}.
\end{aligned}
\end{equation}

\begin{proposition}\label{prop:druga}
 For $\alpha\in\ ]\,\underline{\theta},\overline{\theta}\,[$ \eqref{eq:thetaalpha},
and $s\in\C\sminus\bigcup_{\omega\in 2\pi\i \Z}(\omega+\e^{\i\alpha}\R_{\geq 0})$,
\begin{equation}\label{eq:ThetaQ}
	\Theta_\orbit(s)=\int_{\Circ(\e^{\i\alpha}\R_{\geq 0})} \frac{\tilde Q(\xi,s)}{1-\e^{\,\xi-s}}\d\xi
	=\int_{\e^{\i\alpha}\R_{\geq 0}} \frac{Q(\xi,s)}{1-\e^{\,\xi-s}}\d\xi,
\end{equation}	
where the functions $Q(s,\sigma)$ and $\tilde Q(s,\sigma)$ are defined in \eqref{eq:Q}.
\end{proposition}

\begin{proof}
Using expression \eqref{eq:ThetaC}, and changing the variable of the integration, we get:
\begin{equation}\label{eq:ThetaD}
\begin{aligned}
	\Theta_\orbit(s)&=\int_{b+\i \e^{-\i\alpha}\R}\!(B_m\!\circ\psi_\orbit(x))\,\e^{-s\,t(x)}\d\psi_\orbit(x)\\
	&=\int_{b+\i \e^{-\i\alpha}\R}\!(B_m\!\circ t(x))\,\e^{-s\,t\circ\psi_\orbit^{\circ(-1)}\!\circ t(x)}\d t(x)\\
	&=\int_{b+\i \e^{-\i\alpha}\R}\!(B_m\!\circ t(x))\,\e^{-su(x)}\e^{-s\,t(x)}\d t(x),
\end{aligned}
\end{equation}
for $s$ in the strip
\begin{equation}\label{eq:strip}
\Big\{\tfrac{\im(\e^{-\i \alpha}s)}{\cos\alpha}\in\ ]2\pi(m\!-\!1),\, 2\pi m[\,\Big\}\end{equation} 
between the points $2\pi\i m$ and $2\pi\i(m-1)$, parallel to $\e^{\i\alpha}\R$ (see the discussion after \eqref{eq:ThetaC}).

We want to rewrite the integral \eqref{eq:ThetaD} as a convolution of two Borel transforms. In order to assure convergence, we shall divide the integration line $b+\i\e^{-\i\alpha}\R$ into two rays 
$b+\i\e^{-\i\alpha}\R_{\geq 0}$ and $b-\i \e^{-\i\alpha}\R_{\geq 0}$. 
We denote
\[\tilde\Phi_\pm(s,\sigma):=\int_{b\mp\i \e^{-\i\alpha}\R_{\geq 0}}\!\!\!(B_m\!\circ t(x))\,\e^{-\sigma u(x)}\e^{-s\,t(x)}\d t(x),\]
and $\Phi(s,\sigma):=\tilde\Phi_-(s,\sigma)-\tilde\Phi_+(s,\sigma)=\int_{b+\i\e^{-\i\alpha}\R}(B_m\circ t)\,\e^{-\sigma u}\e^{-st}\d t$, so that $\Theta_\orbit(s)=\Phi(s,s)$.
We rewrite each $\tilde\Phi_\pm(s,\sigma)$ as an integral over the whole line $b+\i\e^{-\i\alpha}\R$
\[\tilde\Phi_\pm(s,\sigma)=\int_{b+\i\e^{-\i\alpha}\R}\!
g_\pm(x,\sigma)\cdot h(x)\,\e^{-s\,t(x)}\d t(x)=
2\pi\i\,\Borel_\alpha\big\{g_\pm\cdot h\big\}(s),
\]
where
$g_\pm(x,\sigma):= (\mathbbm{1}_{b\mp\i\e^{-\i\alpha}\R_{\geq 0} } \circ t(x))\cdot\e^{-\sigma u(x)}$ and $h(x):=B_m\!\circ t(x)$.
Now we apply the standard convolution identity
for two-sided Borel--Laplace--Fourier type transforms of a product of two functions on the line $t(x)\in b+\i\e^{-\i\alpha}\R$: 
\[\tilde \Phi_\pm(s,\sigma)=2\pi\i\,\Borel_\alpha\big\{g_\pm \cdot h\big\}(s)
=2\pi\i\int_{\i\e^{\i\alpha}(0\pms)+\e^{\i\alpha}\R}\!\!\!\!
\Borel_\alpha\big\{g_\pm\big\}(\xi)\cdot\Borel_\alpha\big\{h\big\}(s-\xi)\d\xi,\]
for $s$ in the strip \eqref{eq:strip}.
Here,
\begin{align*}
2\pi\i\,\Borel_\alpha\{h\}(s)&=2\pi\i\,\Borel_\alpha\{B_m\!\circ t(x)\}(s)=
\int_{b+\i\e^{-\i\alpha}\R} \tfrac{\e^{2\pi\i m\,t(x)}}{\e^{2\pi\i\, t(x)}-1}\e^{-s\,t(x)}\d t(x) =\mfrac{1}{1-\e^{-s}}
\end{align*}
by the residue theorem, on the strip \eqref{eq:strip}, and
\[\Borel_\alpha\{g_\pm\}(s)= \tfrac{1}{2\pi\i}\int_{b\mp\i\e^{-\i\alpha}\R_{\geq 0}}\!\!\!\! \e^{-\sigma u(x)}\e^{-s\,t(x)}\d t(x) = \tilde Q(s,\sigma), \]
on the half-plane 
$\big\{\Re(\e^{-\i\beta_\pm}s)>0\big\}$ bounded by the line $\e^{\i\alpha}\R$,
where $\beta_+=\alpha+\frac\pi2$, $\beta_-=\alpha+\frac{3\pi}{2}$. 
We have
\[\Borel_\alpha\{g_-\}(s)-\Borel_\alpha\{g_+\}(s)=\BV[\tilde Q(s,\sigma)]_{\e^{\i\alpha}\R}=
\begin{cases}
\BV[\tilde Q(s,\sigma)]_{\e^{\i\alpha}\R_{\geq0}},\!\!\! & s\in\e^{\i\alpha}\R_{\geq 0},\\0,&s\in\e^{\i\alpha}\R_{<0},
\end{cases}
\]
in distributional sense, hence, on the strip \eqref{eq:strip}:
\begin{equation}\label{eq:PHI}
\Phi(s,\sigma)=\tilde\Phi_-(s,\sigma)-\tilde\Phi_+(s,\sigma)=\int_{\Circ(\e^{\i\alpha}\R_{\geq 0})} \frac{\tilde Q(\xi,\sigma)}{1-\e^{\,\xi-s}}\d\xi.
\end{equation}
Since  the identity \eqref{eq:PHI} is independent of the integer $m\in\Z$ determining the strip \eqref{eq:strip}, it is true on the whole slit plane $s\in\C\sminus\bigcup_{\omega\in 2\pi\i \Z}(\omega+\e^{\i\alpha}\R_{\geq 0})$, and for all $\sigma\in\C$.
The identity \eqref{eq:ThetaQ} now follows by restricting to $\{\sigma=s\}$. 
\end{proof}

\end{document}
